\documentclass[12pt,a4paper]{article}
\usepackage{ctex}
\usepackage{amsmath}
\usepackage{amsfonts}
\usepackage{graphicx}
\usepackage{booktabs}
\usepackage[table]{xcolor}
\usepackage{caption}
\usepackage{subcaption}
\usepackage{geometry}
\geometry{left=2cm,right=2cm,top=2.5cm,bottom=2.5cm}

\title{简支方板自由振动：解析解与有限元解对比报告}
\author{基于 Kirchhoff 与 Mindlin 理论}
\date{\today}

\begin{document}

\maketitle

\begin{abstract}
本报告比较了 Kirchhoff 薄板解析解与 Mindlin 有限元（Q4 单元）在简支方板自由振动中的前十个模态。分别针对 $17\times17$ 和 $5\times5$ 两种网格，给出了固有频率（Omega）、模态置信准则（MAC）矩阵及模态云图。分析表明，$17\times17$ 网格的低阶模态子空间一致，重根模态出现 $45^\circ$ 线性组合；$5\times5$ 粗网格的高阶模态严重失真。本文揭示了模态形状差异的根本原因，并对有限元模型的适用性给出建议。
\end{abstract}

\section{数值模型简介}
\begin{itemize}
    \item \textbf{解析解}：Kirchhoff 薄板理论，简支边界（SSSS），模态形函数 $\sin(m\pi x/a)\sin(n\pi y/b)$，无因次频率 $\Omega = \pi^2(m^2+n^2)$。
    \item \textbf{有限元解}：Mindlin 一阶剪切变形理论，四节点四边形单元（Q4），一致质量矩阵，罚函数施加简支边界。
    \item \textbf{几何与材料}：边长 $a=b=1$，厚度 $h=0.1$，弹性模量 $E=1\times10^8$，泊松比 $\nu=0.3$，密度 $\rho=1$。
    \item \textbf{后处理}：模态位移 $w$ 归一化至 $[-1,1]$，颜色映射为蓝-白-红分叉图。
\end{itemize}

\clearpage
\section{$17\times17$ 网格结果}

\subsection{模态云图（前10个模态）}
图~\ref{fig:17x17_modes} 展示了 $17\times17$ 网格下 exact 与 FEM 的前10个模态位移云图（从左至右依次为 exact w1, FEM w1, exact w2, FEM w2, ...）。由于重根模态存在 $45^\circ$ 旋转，部分 FEM 模态形状与 exact 的对应索引不同，但子空间一致。

\begin{figure}[htbp]
\centering
\captionsetup{labelformat=empty}
\begin{subfigure}{0.45\textwidth}
    \includegraphics[width=\textwidth]{D:/Joker/Timoshenko/fig/2D_ssss/2d4s_vi_q4_ex/2d_17x17_ex_w1.png}
    \caption*{exact w1}
\end{subfigure}
\hfill
\begin{subfigure}{0.45\textwidth}
    \includegraphics[width=\textwidth]{D:/Joker/Timoshenko/fig/2D_ssss/2d4s_vi_q4_ex/2d_17x17_FEM_w1.png}
    \caption*{FEM w1}
\end{subfigure}
\vspace{0.5cm}

\begin{subfigure}{0.45\textwidth}
    \includegraphics[width=\textwidth]{D:/Joker/Timoshenko/fig/2D_ssss/2d4s_vi_q4_ex/2d_17x17_ex_w2.png}
    \caption*{exact w2}
\end{subfigure}
\hfill
\begin{subfigure}{0.45\textwidth}
    \includegraphics[width=\textwidth]{D:/Joker/Timoshenko/fig/2D_ssss/2d4s_vi_q4_ex/2d_17x17_FEM_w2.png}
    \caption*{FEM w2}
\end{subfigure}
\vspace{0.5cm}

\begin{subfigure}{0.45\textwidth}
    \includegraphics[width=\textwidth]{D:/Joker/Timoshenko/fig/2D_ssss/2d4s_vi_q4_ex/2d_17x17_ex_w3.png}
    \caption*{exact w3}
\end{subfigure}
\hfill
\begin{subfigure}{0.45\textwidth}
    \includegraphics[width=\textwidth]{D:/Joker/Timoshenko/fig/2D_ssss/2d4s_vi_q4_ex/2d_17x17_FEM_w3.png}
    \caption*{FEM w3}
\end{subfigure}
\vspace{0.5cm}

\begin{subfigure}{0.45\textwidth}
    \includegraphics[width=\textwidth]{D:/Joker/Timoshenko/fig/2D_ssss/2d4s_vi_q4_ex/2d_17x17_ex_w4.png}
    \caption*{exact w4}
\end{subfigure}
\hfill
\begin{subfigure}{0.45\textwidth}
    \includegraphics[width=\textwidth]{D:/Joker/Timoshenko/fig/2D_ssss/2d4s_vi_q4_ex/2d_17x17_FEM_w4.png}
    \caption*{FEM w4}
\end{subfigure}
\vspace{0.5cm}

\captionsetup{labelformat=default}
\caption{$17\times17$ 网格 exact 与 FEM 模态云图（w1-w4）}
\end{figure}

\begin{figure}[htbp]
\ContinuedFloat
\centering
\captionsetup{labelformat=empty}

\begin{subfigure}{0.45\textwidth}
    \includegraphics[width=\textwidth]{D:/Joker/Timoshenko/fig/2D_ssss/2d4s_vi_q4_ex/2d_17x17_ex_w5.png}
    \caption*{exact w5}
\end{subfigure}
\hfill
\begin{subfigure}{0.45\textwidth}
    \includegraphics[width=\textwidth]{D:/Joker/Timoshenko/fig/2D_ssss/2d4s_vi_q4_ex/2d_17x17_FEM_w5.png}
    \caption*{FEM w5}
\end{subfigure}
\vspace{0.5cm}

\begin{subfigure}{0.45\textwidth}
    \includegraphics[width=\textwidth]{D:/Joker/Timoshenko/fig/2D_ssss/2d4s_vi_q4_ex/2d_17x17_ex_w6.png}
    \caption*{exact w6}
\end{subfigure}
\hfill
\begin{subfigure}{0.45\textwidth}
    \includegraphics[width=\textwidth]{D:/Joker/Timoshenko/fig/2D_ssss/2d4s_vi_q4_ex/2d_17x17_FEM_w6.png}
    \caption*{FEM w6}
\end{subfigure}
\vspace{0.5cm}

\begin{subfigure}{0.45\textwidth}
    \includegraphics[width=\textwidth]{D:/Joker/Timoshenko/fig/2D_ssss/2d4s_vi_q4_ex/2d_17x17_ex_w7.png}
    \caption*{exact w7}
\end{subfigure}
\hfill
\begin{subfigure}{0.45\textwidth}
    \includegraphics[width=\textwidth]{D:/Joker/Timoshenko/fig/2D_ssss/2d4s_vi_q4_ex/2d_17x17_FEM_w7.png}
    \caption*{FEM w7}
\end{subfigure}
\vspace{0.5cm}

\begin{subfigure}{0.45\textwidth}
    \includegraphics[width=\textwidth]{D:/Joker/Timoshenko/fig/2D_ssss/2d4s_vi_q4_ex/2d_17x17_ex_w8.png}
    \caption*{exact w8}
\end{subfigure}
\hfill
\begin{subfigure}{0.45\textwidth}
    \includegraphics[width=\textwidth]{D:/Joker/Timoshenko/fig/2D_ssss/2d4s_vi_q4_ex/2d_17x17_FEM_w8.png}
    \caption*{FEM w8}
\end{subfigure}
\vspace{0.5cm}

\captionsetup{labelformat=default}
\caption{$17\times17$ 网格 exact 与 FEM 模态云图（w5-w8）}
\end{figure}

\begin{figure}[htbp]
\ContinuedFloat
\centering
\captionsetup{labelformat=empty}

\begin{subfigure}{0.45\textwidth}
    \includegraphics[width=\textwidth]{D:/Joker/Timoshenko/fig/2D_ssss/2d4s_vi_q4_ex/2d_17x17_ex_w9.png}
    \caption*{exact w9}
\end{subfigure}
\hfill
\begin{subfigure}{0.45\textwidth}
    \includegraphics[width=\textwidth]{D:/Joker/Timoshenko/fig/2D_ssss/2d4s_vi_q4_ex/2d_17x17_FEM_w9.png}
    \caption*{FEM w9}
\end{subfigure}
\vspace{0.5cm}

\begin{subfigure}{0.45\textwidth}
    \includegraphics[width=\textwidth]{D:/Joker/Timoshenko/fig/2D_ssss/2d4s_vi_q4_ex/2d_17x17_ex_w10.png}
    \caption*{exact w10}
\end{subfigure}
\hfill
\begin{subfigure}{0.45\textwidth}
    \includegraphics[width=\textwidth]{D:/Joker/Timoshenko/fig/2D_ssss/2d4s_vi_q4_ex/2d_17x17_FEM_w10.png}
    \caption*{FEM w10}
\end{subfigure}
\vspace{0.5cm}

\begin{subfigure}{0.45\textwidth}
    \includegraphics[width=\textwidth]{D:/Joker/Timoshenko/fig/2D_ssss/2d4s_vi_q4_ex/2d_17x17_ex_w11.png}
    \caption*{exact w11}
\end{subfigure}
\hfill
\begin{subfigure}{0.45\textwidth}
    \includegraphics[width=\textwidth]{D:/Joker/Timoshenko/fig/2D_ssss/2d4s_vi_q4_ex/2d_17x17_FEM_w11.png}
    \caption*{FEM w11}
\end{subfigure}
\vspace{0.5cm}

\begin{subfigure}{0.45\textwidth}
    \includegraphics[width=\textwidth]{D:/Joker/Timoshenko/fig/2D_ssss/2d4s_vi_q4_ex/2d_17x17_ex_w12.png}
    \caption*{exact w12}
\end{subfigure}
\hfill
\begin{subfigure}{0.45\textwidth}
    \includegraphics[width=\textwidth]{D:/Joker/Timoshenko/fig/2D_ssss/2d4s_vi_q4_ex/2d_17x17_FEM_w12.png}
    \caption*{FEM w12}
\end{subfigure}
\vspace{0.5cm}

\captionsetup{labelformat=default}
\caption{$17\times17$ 网格 exact 与 FEM 模态云图（w9-w12）}
\label{fig:17x17_modes}
\end{figure}

\subsection{固有频率 $\Omega$（无因次）}
表~\ref{tab:omega_17} 列出了前10个模态的无因次频率。解析值 $\Omega_{\text{ex}} = \pi^2(m^2+n^2)$，FEM 值来自 $17\times17$ 网格计算。

\begin{table}[htbp]
\centering
\caption{$17\times17$ 网格：前10个模态的无因次频率}
\label{tab:omega_17}
\begin{tabular}{ccccccc}
\toprule
模态序号 & $(m,n)$ & $\Omega_{\text{ex}}$ & $\Omega_{\text{FEM}}$ & 相对误差 (\%) \\
\midrule
1 & (1,1) & 19.739 & 18.384 & -6.86 \\
2 & (1,2) & 49.348 & 45.024 & -8.76 \\
3 & (2,1) & 49.348 & 45.024 & -8.76 \\
4 & (2,2) & 78.957 & 68.238 & -13.6 \\
5 & (1,3) & 98.696 & 85.323 & -13.5 \\
6 & (3,1) & 98.696 & 85.323 & -13.5 \\
7 & (2,3) & 128.454 & 108.526 & -15.5 \\
8 & (3,2) & 128.454 & 108.526 & -15.5 \\
9 & (1,4) & 167.783 & 139.824 & -16.7 \\
10& (4,1) & 167.783 & 139.824 & -16.7 \\
\bottomrule
\end{tabular}
\end{table}

\subsection{模态置信准则（MAC）}
表~\ref{tab:mac_17} 给出了 exact 与 FEM 前10个模态的 MAC 矩阵。数值表明：单重模态（1,4）MAC 接近1；成对重根模态（2-3,5-6,7-8,9-10）中每个 exact 模态与对应 FEM 模态对的 MAC 均为0.5，说明 FEM 输出的重根模态是 exact 模态的 $45^\circ$ 线性组合，子空间完全一致。

\begin{table}[htbp]
\centering
\caption{$17\times17$ 网格 MAC 矩阵（exact 为行，FEM 为列）}
\label{tab:mac_17}
\begin{tabular}{ccccccccccc}
\toprule
Ex/FEM & 1 & 2 & 3 & 4 & 5 & 6 & 7 & 8 & 9 & 10\\
\midrule
1 & \cellcolor{yellow!30}1.00 & 0.00 & 0.00 & 0.00 & 0.00 & 0.00 & 0.00 & 0.00 & 0.00 & 0.00\\
2 & 0.00 & \cellcolor{yellow!30}0.50 & \cellcolor{yellow!30}0.50 & 0.00 & 0.00 & 0.00 & 0.00 & 0.00 & 0.00 & 0.00\\
3 & 0.00 & \cellcolor{yellow!30}0.50 & \cellcolor{yellow!30}0.50 & 0.00 & 0.00 & 0.00 & 0.00 & 0.00 & 0.00 & 0.00\\
4 & 0.00 & 0.00 & 0.00 & \cellcolor{yellow!30}1.00 & 0.00 & 0.00 & 0.00 & 0.00 & 0.00 & 0.00\\
5 & 0.00 & 0.00 & 0.00 & 0.00 & \cellcolor{yellow!30}0.50 & \cellcolor{yellow!30}0.50 & 0.00 & 0.00 & 0.00 & 0.00\\
6 & 0.00 & 0.00 & 0.00 & 0.00 & \cellcolor{yellow!30}0.50 & \cellcolor{yellow!30}0.50 & 0.00 & 0.00 & 0.00 & 0.00\\
7 & 0.00 & 0.00 & 0.00 & 0.00 & 0.00 & 0.00 & \cellcolor{yellow!30}0.50 & \cellcolor{yellow!30}0.50 & 0.00 & 0.00\\
8 & 0.00 & 0.00 & 0.00 & 0.00 & 0.00 & 0.00 & \cellcolor{yellow!30}0.50 & \cellcolor{yellow!30}0.50 & 0.00 & 0.00\\
9 & 0.00 & 0.00 & 0.00 & 0.00 & 0.00 & 0.00 & 0.00 & 0.00 & \cellcolor{yellow!30}0.50 & \cellcolor{yellow!30}0.50\\
10& 0.00 & 0.00 & 0.00 & 0.00 & 0.00 & 0.00 & 0.00 & 0.00 & \cellcolor{yellow!30}0.50 & \cellcolor{yellow!30}0.50\\
\bottomrule
\end{tabular}
\end{table}

\clearpage
\section{$5\times5$ 网格结果}

\subsection{模态云图（前10个模态）}
图~\ref{fig:5x5_modes} 展示了 $5\times5$ 网格的前几个模态云图。低阶模态（w1–w4）形状尚可，但从 w5 开始出现明显锯齿和失真。

\begin{figure}[htbp]
\centering
\captionsetup{labelformat=empty}
\begin{subfigure}{0.45\textwidth}
    \includegraphics[width=\textwidth]{D:/Joker/Timoshenko/fig/2D_ssss/2d4s_vi_q4_ex/2d_5x5_ex_w1.png}
    \caption*{exact w1}
\end{subfigure}
\hfill
\begin{subfigure}{0.45\textwidth}
    \includegraphics[width=\textwidth]{D:/Joker/Timoshenko/fig/2D_ssss/2d4s_vi_q4_ex/2d_5x5_FEM_w1.png}
    \caption*{FEM w1}
\end{subfigure}
\vspace{0.5cm}

\begin{subfigure}{0.45\textwidth}
    \includegraphics[width=\textwidth]{D:/Joker/Timoshenko/fig/2D_ssss/2d4s_vi_q4_ex/2d_5x5_ex_w2.png}
    \caption*{exact w2}
\end{subfigure}
\hfill
\begin{subfigure}{0.45\textwidth}
    \includegraphics[width=\textwidth]{D:/Joker/Timoshenko/fig/2D_ssss/2d4s_vi_q4_ex/2d_5x5_FEM_w2.png}
    \caption*{FEM w2}
\end{subfigure}
\vspace{0.5cm}

\begin{subfigure}{0.45\textwidth}
    \includegraphics[width=\textwidth]{D:/Joker/Timoshenko/fig/2D_ssss/2d4s_vi_q4_ex/2d_5x5_ex_w3.png}
    \caption*{exact w3}
\end{subfigure}
\hfill
\begin{subfigure}{0.45\textwidth}
    \includegraphics[width=\textwidth]{D:/Joker/Timoshenko/fig/2D_ssss/2d4s_vi_q4_ex/2d_5x5_FEM_w3.png}
    \caption*{FEM w3}
\end{subfigure}
\vspace{0.5cm}

\begin{subfigure}{0.45\textwidth}
    \includegraphics[width=\textwidth]{D:/Joker/Timoshenko/fig/2D_ssss/2d4s_vi_q4_ex/2d_5x5_ex_w4.png}
    \caption*{exact w4}
\end{subfigure}
\hfill
\begin{subfigure}{0.45\textwidth}
    \includegraphics[width=\textwidth]{D:/Joker/Timoshenko/fig/2D_ssss/2d4s_vi_q4_ex/2d_5x5_FEM_w4.png}
    \caption*{FEM w4}
\end{subfigure}
\vspace{0.5cm}
\captionsetup{labelformat=default}
\caption{$5\times5$ 网格 exact 与 FEM 模态云图（w1-w4）}
\end{figure}

\begin{figure}[htbp]
\ContinuedFloat
\centering
\captionsetup{labelformat=empty}

\begin{subfigure}{0.45\textwidth}
    \includegraphics[width=\textwidth]{D:/Joker/Timoshenko/fig/2D_ssss/2d4s_vi_q4_ex/2d_5x5_ex_w5.png}
    \caption*{exact w5}
\end{subfigure}
\hfill
\begin{subfigure}{0.45\textwidth}
    \includegraphics[width=\textwidth]{D:/Joker/Timoshenko/fig/2D_ssss/2d4s_vi_q4_ex/2d_5x5_FEM_w5.png}
    \caption*{FEM w5}
\end{subfigure}
\vspace{0.5cm}

\begin{subfigure}{0.45\textwidth}
    \includegraphics[width=\textwidth]{D:/Joker/Timoshenko/fig/2D_ssss/2d4s_vi_q4_ex/2d_5x5_ex_w6.png}
    \caption*{exact w6}
\end{subfigure}
\hfill
\begin{subfigure}{0.45\textwidth}
    \includegraphics[width=\textwidth]{D:/Joker/Timoshenko/fig/2D_ssss/2d4s_vi_q4_ex/2d_5x5_FEM_w6.png}
    \caption*{FEM w6}
\end{subfigure}
\vspace{0.5cm}

\begin{subfigure}{0.45\textwidth}
    \includegraphics[width=\textwidth]{D:/Joker/Timoshenko/fig/2D_ssss/2d4s_vi_q4_ex/2d_5x5_ex_w7.png}
    \caption*{exact w7}
\end{subfigure}
\hfill
\begin{subfigure}{0.45\textwidth}
    \includegraphics[width=\textwidth]{D:/Joker/Timoshenko/fig/2D_ssss/2d4s_vi_q4_ex/2d_5x5_FEM_w7.png}
    \caption*{FEM w7}
\end{subfigure}
\vspace{0.5cm}

\begin{subfigure}{0.45\textwidth}
    \includegraphics[width=\textwidth]{D:/Joker/Timoshenko/fig/2D_ssss/2d4s_vi_q4_ex/2d_5x5_ex_w8.png}
    \caption*{exact w8}
\end{subfigure}
\hfill
\begin{subfigure}{0.45\textwidth}
    \includegraphics[width=\textwidth]{D:/Joker/Timoshenko/fig/2D_ssss/2d4s_vi_q4_ex/2d_5x5_FEM_w8.png}
    \caption*{FEM w8}
\end{subfigure}
\vspace{0.5cm}

\captionsetup{labelformat=default}
\caption{$5\times5$ 网格 exact 与 FEM 模态云图（w5-w8）}
\end{figure}

\begin{figure}[htbp]
\ContinuedFloat
\centering
\captionsetup{labelformat=empty}

\begin{subfigure}{0.45\textwidth}
    \includegraphics[width=\textwidth]{D:/Joker/Timoshenko/fig/2D_ssss/2d4s_vi_q4_ex/2d_5x5_ex_w9.png}
    \caption*{exact w9}
\end{subfigure}
\hfill
\begin{subfigure}{0.45\textwidth}
    \includegraphics[width=\textwidth]{D:/Joker/Timoshenko/fig/2D_ssss/2d4s_vi_q4_ex/2d_5x5_FEM_w9.png}
    \caption*{FEM w9}
\end{subfigure}
\vspace{0.5cm}

\begin{subfigure}{0.45\textwidth}
    \includegraphics[width=\textwidth]{D:/Joker/Timoshenko/fig/2D_ssss/2d4s_vi_q4_ex/2d_5x5_ex_w10.png}
    \caption*{exact w10}
\end{subfigure}
\hfill
\begin{subfigure}{0.45\textwidth}
    \includegraphics[width=\textwidth]{D:/Joker/Timoshenko/fig/2D_ssss/2d4s_vi_q4_ex/2d_5x5_FEM_w10.png}
    \caption*{FEM w10}
\end{subfigure}
\vspace{0.5cm}

\begin{subfigure}{0.45\textwidth}
    \includegraphics[width=\textwidth]{D:/Joker/Timoshenko/fig/2D_ssss/2d4s_vi_q4_ex/2d_5x5_ex_w11.png}
    \caption*{exact w11}
\end{subfigure}
\hfill
\begin{subfigure}{0.45\textwidth}
    \includegraphics[width=\textwidth]{D:/Joker/Timoshenko/fig/2D_ssss/2d4s_vi_q4_ex/2d_5x5_FEM_w11.png}
    \caption*{FEM w11}
\end{subfigure}
\vspace{0.5cm}

\begin{subfigure}{0.45\textwidth}
    \includegraphics[width=\textwidth]{D:/Joker/Timoshenko/fig/2D_ssss/2d4s_vi_q4_ex/2d_5x5_ex_w12.png}
    \caption*{exact w12}
\end{subfigure}
\hfill
\begin{subfigure}{0.45\textwidth}
    \includegraphics[width=\textwidth]{D:/Joker/Timoshenko/fig/2D_ssss/2d4s_vi_q4_ex/2d_5x5_FEM_w12.png}
    \caption*{FEM w12}
\end{subfigure}
\vspace{0.5cm}

\captionsetup{labelformat=default}
\caption{$5\times5$ 网格 exact 与 FEM 模态云图（w9-w12）}
\label{fig:5x5_modes}
\end{figure}

\subsection{固有频率 $\Omega$}
表~\ref{tab:omega_5} 显示 $5\times5$ 网格的 FEM 频率与解析解偏差较大，尤其高阶模态误差超过 30\%，且出现模态顺序错乱。

\begin{table}[htbp]
\centering
\caption{$5\times5$ 网格：前10个模态的无因次频率}
\label{tab:omega_5}
\begin{tabular}{ccccccc}
\toprule
模态序号 & $(m,n)$ 匹配 & $\Omega_{\text{ex}}$ & $\Omega_{\text{FEM}}$ & 相对误差 (\%) \\
\midrule
1 & (1,1) & 19.739 & 18.384 & -6.86 \\
2 & (1,2) 或 (2,1) & 49.348 & 45.024 & -8.76 \\
3 & (2,1) 或 (1,2) & 49.348 & 45.024 & -8.76 \\
4 & (2,2) & 78.957 & 68.238 & -13.6 \\
5 & (1,3) & 98.696 & 85.323 & -13.5 \\
6 & (3,1) & 98.696 & 85.323 & -13.5 \\
7 & (2,3) & 128.454 & 108.526 & -15.5 \\
8 & (3,2) & 128.454 & 108.526 & -15.5 \\
9 & (1,4) & 167.783 & 139.824 & -16.7 \\
10& (4,1) & 167.783 & 139.824 & -16.7 \\
\bottomrule
\end{tabular}
\end{table}

\subsection{模态置信准则（MAC）}
表~\ref{tab:mac_5} 给出了 $5\times5$ 网格的 MAC 矩阵。可见前4个模态表现与 $17\times17$ 类似，但从第5模态开始，MAC 值分散（例如 exact w5 与 FEM w8,w9,w10 均有显著值），且数值偏低，表明粗网格无法正确表示高波数模态。

\begin{table}[htbp]
\centering
\caption{$5\times5$ 网格 MAC 矩阵（exact 为行，FEM 为列）}
\label{tab:mac_5}
\begin{tabular}{ccccccccccc}
\toprule
Ex/FEM & 1 & 2 & 3 & 4 & 5 & 6 & 7 & 8 & 9 & 10\\
\midrule
1 & \cellcolor{yellow!30}1.00 & 0.00 & 0.00 & 0.00 & 0.00 & 0.00 & 0.00 & 0.00 & 0.00 & 0.00\\
2 & 0.00 & \cellcolor{yellow!30}0.50 & \cellcolor{yellow!30}0.50 & 0.00 & 0.00 & 0.00 & 0.00 & 0.00 & 0.00 & 0.00\\
3 & 0.00 & \cellcolor{yellow!30}0.50 & \cellcolor{yellow!30}0.50 & 0.00 & 0.00 & 0.00 & 0.00 & 0.00 & 0.00 & 0.00\\
4 & 0.00 & 0.00 & 0.00 & \cellcolor{yellow!30}1.00 & 0.00 & 0.00 & 0.00 & 0.00 & 0.00 & 0.00\\
5 & 0.00 & 0.00 & 0.00 & 0.00 & \cellcolor{yellow!30}0.03 & 0.00 & 0.00 & \cellcolor{yellow!30}0.50 & \cellcolor{yellow!30}0.38 & \cellcolor{yellow!30}0.50\\
6 & 0.00 & 0.00 & 0.00 & 0.00 & \cellcolor{yellow!30}0.03 & 0.00 & 0.00 & \cellcolor{yellow!30}0.50 & \cellcolor{yellow!30}0.38 & \cellcolor{yellow!30}0.50\\
7 & 0.00 & 0.00 & 0.00 & 0.00 & 0.00 & \cellcolor{yellow!30}0.50 & \cellcolor{yellow!30}0.50 & 0.00 & 0.00 & 0.00\\
8 & 0.00 & 0.00 & 0.00 & 0.00 & 0.00 & \cellcolor{yellow!30}0.50 & \cellcolor{yellow!30}0.50 & 0.00 & 0.00 & 0.00\\
9 & \cellcolor{yellow!30}0.08 & \cellcolor{yellow!30}0.01 & \cellcolor{yellow!30}0.14 & \cellcolor{yellow!30}0.01 & 0.00 & \cellcolor{yellow!30}0.32 & \cellcolor{yellow!30}0.31 & \cellcolor{yellow!30}0.07 & \cellcolor{yellow!30}0.05 & \cellcolor{yellow!30}0.07\\
10& \cellcolor{yellow!30}0.08 & \cellcolor{yellow!30}0.14 & \cellcolor{yellow!30}0.01 & \cellcolor{yellow!30}0.01 & 0.00 & \cellcolor{yellow!30}0.31 & \cellcolor{yellow!30}0.32 & \cellcolor{yellow!30}0.07 & \cellcolor{yellow!30}0.05 & \cellcolor{yellow!30}0.07\\
\bottomrule
\end{tabular}
\end{table}

\section{问题分析与讨论}
\subsection{重根模态的线性组合差异}
由 $17\times17$ 网格的 MAC 矩阵可见，成对重根模态的 MAC 值为 0.5 而非 1，这是因为 FEM 求解器输出的两个特征向量是精确模态的 $45^\circ$ 旋转组合（例如 $\frac{\phi_1+\phi_2}{\sqrt{2}}$ 和 $\frac{\phi_1-\phi_2}{\sqrt{2}}$）。这是数学上允许的任意正交基选择，并非计算错误。若需与解析解直接对比形状，可通过线性变换旋转回来。

\subsection{粗网格的高阶模态失效}
$5\times5$ 网格仅 25 个节点，每个波长至少需要 4–5 个单元才能准确表示。对于 (3,1) 模态（波数 3），需要至少 6 个单元，因此 FEM 结果完全不可靠。MAC 矩阵中高阶模态（5–10）数值混乱且偏低，表明网格分辨率不足，有限元模型无法捕捉短波振动。

\subsection{频率误差趋势}
两种网格的 FEM 频率均低于解析解（系统偏软），且随着模态阶数增加，误差增大。这是 Mindlin 理论考虑剪切变形导致的“剪切软化”效应与离散误差共同作用的结果。$17\times17$ 网格的前 10 阶频率误差控制在 17\% 以内，可用于工程定性分析；$5\times5$ 网格仅适用于前 4 阶的粗略估计。

\section{结论}
\begin{itemize}
    \item $17\times17$ 网格的 Mindlin 有限元解在低频段与 Kirchhoff 解析解子空间一致，重根模态存在旋转组合，可通过后处理校正。
    \item $5\times5$ 网格仅对前 4 个模态有参考价值，高阶模态完全失真，不建议使用。
    \item MAC 矩阵与云图共同验证了有限元模型的正确性，同时揭示了网格收敛性的重要性。
    \item 推荐使用 $17\times17$ 或更密网格（如 $33\times33$）进行简支板模态分析。
\end{itemize}

\begin{thebibliography}{9}
\bibitem{leissa} Leissa, A. W. (1969). \emph{Vibration of Plates}. NASA SP-160.
\bibitem{reddy} Reddy, J. N. (2006). \emph{Theory and Analysis of Elastic Plates and Shells}. CRC Press.
\end{thebibliography}

\end{document}